% ════════════════════════════════════════════
% 具体分析（5.2.1）
% ════════════════════════════════════════════
\subsubsection{整理思路}

本部分属于数据清洗与特征整理，采用规则核查、众数填补、独热编码与对数变换相结合的方法。清洗的目标不是把数据“修得好看”，而是让每一列都有明确含义、可追溯的处理记录，并为回归提供数值稳定的设计矩阵。

首先逐列检查缺失、重复与越界，确认唯一的缺失发生在付费标记上，并对曝光量的极端值作业务真实性判断。接着在缺失机制假设下选择众数填补，同时保留删除缺失行的对照方案，以便量化填补对结果的影响。然后按变量类型完成编码：名义变量采用独热编码并设定参照水平，周期性时间变量保留哑变量，发帖小时因分布过于集中而合并为三个时段组。最后对曝光量做 log1p 变换、对粉丝规模做对数变换，比较变换前后的偏度与峰度，并把所有处理结果写入统一分析数据。具体分析流程如图\ref{fig:analysis2_flow}所示。

\begin{figure}[ht]
    \begin{center}
        \begin{tikzpicture}[x=1cm, y=0.7cm]
            \node[box] (n11) at (0, 0) {逐列缺失检查};
            \node[box] (n12) at (5, 0) {异常值核查};
            \node[box] (n13) at (10, 0) {重复记录检查};
            \node[box] (n22) at (5, -3) {变量类型编码};
            \node[box] (n21) at (0, -3) {缺失值填补};
            \node[box] (n23) at (10, -3) {稀疏时段合并};
            \node[box] (n31) at (0, -6) {对数变换};
            \node[box] (n32) at (5, -6) {分布前后对比};
            \node[box] (n33) at (10, -6) {敏感性检验};
            \draw[arrow] (n11) -- (n12);
            \draw[arrow] (n12) -- (n13);
            \draw[arrow] (n23) -- (n22);
            \draw[arrow] (n22) -- (n21);
            \draw[arrow] (n13) -- ++(0,-1.0) -| (n23);
            \draw[arrow] (n21) -- (n31);
            \draw[arrow] (n31) -- (n32);
            \draw[arrow] (n32) -- (n33);
        \end{tikzpicture}
        \caption{数据整理流程图}
        \label{fig:analysis2_flow}
    \end{center}
\end{figure}
